% !TEX root = ../algebra_02.tex

\section{Характеристика внутренней прямой суммы в терминах базиса}

Понятие базиса пространства тесно связано с прямой суммой одномерных подпространств.
\begin{Remark}{}{}
	Пусть $e_1, \ldots, e_k$ — ненулевые элементы пространства $V$. Тогда $e_1, \ldots, e_k$ образуют базис $V$ тогда и только тогда, когда пространство раскладывается в прямую сумму их одномерных линейных оболочек:
	\[
	V = \mathrm{Lin}(e_1) \oplus \cdots \oplus \mathrm{Lin}(e_k).
	\]
\end{Remark}

Для произвольных подпространств прямая сумма также напрямую связана с объединением их базисов.

\begin{Theorem}{Базис прямой суммы}{}
	Пусть $W_1, \ldots, W_k \leq V$, и пусть $e_{i,1}, \ldots, e_{i,d_i}$ — базис $W_i$. Тогда
	\[
	V = W_1 \oplus \cdots \oplus W_k \iff \{e_{i,j} \mid 1 \leq i \leq k,\; 1 \leq j \leq d_i\} \text{ — базис } V.
	\]
\end{Theorem}

\begin{proof}
	$(\Rightarrow)$: Любой $v \in V$ единственным образом раскладывается как $v = \sum_i w_i$, $w_i \in W_i$, а каждое $w_i$ единственным образом раскладывается по базису $W_i$. Таким образом, $\{e_{i,j}\}$ — базис $V$.
	
	$(\Leftarrow)$: Из того, что $\{e_{i,j}\}$ — базис $V$, следует, что любой вектор единственным образом раскладывается по ним, что равносильно единственности разложения $v = \sum_i w_i$ с $w_i \in W_i$.
\end{proof}